\section{总结与延伸}

\subsection{本讲核心线索}

本讲从最基本的感知机出发，沿着一条清晰的演进路线逐步构建了对序列建模的完整理解：

\begin{enumerate}
    \item \textbf{感知机/前馈网络}：能处理单个时间步的输入，但无法利用历史信息
    \item \textbf{RNN}：通过隐藏状态引入``记忆''，实现了时间步之间的信息传递
    \item \textbf{LSTM}：通过门控机制缓解梯度消失问题，改善长期依赖的学习
    \item \textbf{Attention/Transformer}：完全放弃循环结构，通过自注意力机制直接建模全局依赖关系
\end{enumerate}

\begin{importantbox}{从 RNN 到 Transformer 的范式转变}
这一演进路线体现了深度学习中一个重要的设计哲学转变：
\begin{itemize}
    \item \textbf{RNN 哲学}：按时间顺序处理，维护一个压缩的记忆状态
    \item \textbf{Transformer 哲学}：一次看到全局，让网络自己决定关注什么
\end{itemize}
Transformer 的成功表明：与其让人类设计信息流动的方式（如循环），不如让网络通过学习自行发现数据中的重要模式。
\end{importantbox}

\subsection{关键概念速查}

\begin{table}[H]
\centering
\caption{本讲关键概念速查表}
\begin{tabular}{ll}
\toprule
\textbf{概念} & \textbf{核心要点} \\
\midrule
Hidden State ($h_t$) & RNN 在时间步之间传递的内部记忆 \\
Recurrence Relation & $h_t = f_W(x_t, h_{t-1})$，循环更新隐藏状态 \\
BPTT & 沿时间轴反向传播梯度的算法 \\
Vanishing Gradient & 梯度随时间步衰减，阻碍长期依赖学习 \\
LSTM & 通过门控机制保护梯度流的 RNN 变体 \\
Embedding & 将离散符号映射为连续向量的可学习表示 \\
Positional Encoding & 为 Transformer 输入注入位置信息 \\
Self-Attention & 序列内部的 Query-Key-Value 注意力计算 \\
Scaled Dot-Product & $\text{softmax}(QK^T/\sqrt{d_k})V$ \\
Multi-Head Attention & 多组并行 Attention，捕捉多种关注模式 \\
Transformer & 完全基于 Attention 的序列模型架构 \\
\bottomrule
\end{tabular}
\end{table}

\subsection{拓展阅读}

\begin{itemize}
    \item Vaswani et al., \textit{Attention Is All You Need}, NeurIPS 2017. Transformer 的开创性论文。
    \item Hochreiter \& Schmidhuber, \textit{Long Short-Term Memory}, Neural Computation 1997. LSTM 的原始论文。
    \item Devlin et al., \textit{BERT: Pre-training of Deep Bidirectional Transformers for Language Understanding}, NAACL 2019.
    \item Brown et al., \textit{Language Models are Few-Shot Learners (GPT-3)}, NeurIPS 2020.
    \item Dosovitskiy et al., \textit{An Image is Worth 16x16 Words: Transformers for Image Recognition at Scale (ViT)}, ICLR 2021.
    \item MIT 6.S191 课程官网：\url{http://introtodeeplearning.com}
    \item MIT 6.S191 Lab 1：Deep Learning in Python and Music Generation with RNNs
\end{itemize}

%% --- Fin del contenido --- %%

\end{document}
